\section{Discussion}
\label{sec:discussion}

\revised{FDIL rests on the two ambiguities of intent recognition. Keeping them separate makes each function independently justified, measurable, and deployable. It suits offline or high-throughput visual-content analysis, since its external semantic knowledge is prepared offline, the rationale generator and text teacher are discarded after training, and SLR-C adds only lightweight reranking and residual refinement on top of the shared frozen encoder, which remains the dominant inference cost.}

The limitations lie in the prior and the visual representation. Offline text generation and prior construction add preparation cost, gains shrink on weaker backbones such as ResNet101, and richer priors are not automatically better unless their interaction with the student is controlled. For a few broad, overlapping categories, the residual student can resolve the candidate set toward a competing neighbor and lose accuracy on the original class. \revised{Generalization beyond the intent taxonomy is also limited. Across the transfer studies, the component-level tendencies are broadly consistent, with the semantic prior moving ranking and pointwise distribution fit and the uncertainty-guided teacher moving thresholded F1 and rank correlation, although the overall gains are small and metric-dependent and no single metric improves everywhere.} Natural next steps are source-aware prior weighting instead of uniform averaging, stronger students for weaker backbones, and confusion-pair-specific textual evidence.
